\chapter{Manual de usuario de TierraViva}
\label{anexo:manual_usuario}

Este anexo describe el uso del sistema TierraViva desde el punto de vista del usuario final. A diferencia del manual de instalación, centrado en el despliegue técnico, este manual explica cómo acceder al \textit{dashboard}, cómo interpretar las lecturas, cómo comprobar el estado de las estaciones y cómo actuar ante situaciones habituales de funcionamiento.

\section{Alcance del manual}

Este manual cubre las operaciones de uso habituales:

\begin{itemize}
    \item acceso al \textit{dashboard} web;
    \item selección de estación;
    \item interpretación de lecturas actuales;
    \item interpretación del estado de humedad del suelo;
    \item revisión del estado de conexión;
    \item consulta de gráficas;
    \item revisión del resumen histórico;
    \item consulta de la tabla de últimas lecturas;
    \item interpretación de eventos y códigos de diagnóstico;
    \item comprobaciones básicas ante fallos.
\end{itemize}

No se incluyen en este manual las tareas de montaje eléctrico, instalación en Raspberry Pi o configuración de ThingSpeak, ya que dichas tareas se documentan en el Anexo \ref{anexo:manual_instalacion_sistema}. Tampoco se documentan en detalle los endpoints de la API, que se recogen en el Anexo \ref{anexo:manual_api}.

\section{Perfil de usuario}

El \textit{dashboard} de TierraViva puede ser utilizado por dos tipos principales de usuario:

\begin{table}[H]
\centering
\renewcommand{\arraystretch}{1.25}
\small
\begin{tabular}{|p{0.28\textwidth}|p{0.58\textwidth}|}
\hline
\rowcolor{gray!20}
\textbf{Perfil} & \textbf{Uso principal del sistema} \\
\hline
Usuario no técnico & Consultar el estado de la parcela, comprobar si el suelo está seco o húmedo y revisar si la estación está actualizando datos. \\
\hline
Usuario técnico & Verificar lecturas, revisar estado de conexión, comprobar eventos de sistema, analizar histórico y detectar posibles problemas de sensores, comunicación o \textit{backend}. \\
\hline
\end{tabular}
\caption{Perfiles de usuario previstos para TierraViva.}
\label{tab:perfiles_usuario_tierraviva}
\end{table}

\section{Acceso al \textit{dashboard}}

Una vez arrancado el sistema en la Raspberry Pi, el usuario accede al \textit{dashboard} desde un navegador web conectado a la misma red local.

\begin{verbatim}
http://<IP_RASPBERRY>:8000/
\end{verbatim}

Por ejemplo:

\begin{verbatim}
http://192.168.1.50:8000/
\end{verbatim}

La dirección exacta depende de la IP asignada a la Raspberry Pi en la red local. Si se accede desde la propia Raspberry Pi, también puede utilizarse:

\begin{verbatim}
http://127.0.0.1:8000/
\end{verbatim}

\section{Vista general del \textit{dashboard}}

El \textit{dashboard} está organizado en bloques visuales para que el estado del sistema pueda interpretarse rápidamente. La Tabla \ref{tab:zonas_dashboard} resume las zonas principales de la interfaz.

\begin{table}[H]
\centering
\renewcommand{\arraystretch}{1.25}
\small
\begin{tabular}{|p{0.30\textwidth}|p{0.55\textwidth}|}
\hline
\rowcolor{gray!20}
\textbf{Zona del dashboard} & \textbf{Información mostrada} \\
\hline
Cabecera & Nombre del sistema, estado de conexión y selector de estación. \\
\hline
Tarjetas de lectura actual & Humedad del suelo, temperatura, humedad ambiental, luminosidad y presión. \\
\hline
Estado de última lectura & Frescura del dato, hora de actualización y estado general. \\
\hline
Estado/evento & Código o evento de funcionamiento recibido desde la estación. \\
\hline
Debug/código & Código auxiliar de diagnóstico. \\
\hline
Gráficas & Evolución reciente de las variables almacenadas. \\
\hline
Resumen histórico & Estadísticas de las últimas 24 horas. \\
\hline
Tabla de últimas lecturas & Registros recientes almacenados en SQLite. \\
\hline
\end{tabular}
\caption{Zonas principales del \textit{dashboard} TierraViva.}
\label{tab:zonas_dashboard}
\end{table}

\section{Selección de estación}

Si el sistema tiene más de una estación configurada, el \textit{dashboard} permite seleccionar la estación activa mediante un desplegable. En la configuración actual están disponibleslas estaciones \texttt{A} y \texttt{B}.

Al seleccionar una estación, el \textit{dashboard} actualiza:

\begin{itemize}
    \item lectura actual;
    \item estado de humedad;
    \item gráficas;
    \item resumen histórico;
    \item tabla de últimas lecturas;
    \item códigos de evento y diagnóstico.
\end{itemize}

Esta separación evita mezclar datos de distintas estaciones y permite comparar zonas diferentes de una parcela o distintos entornos de prueba.

\section{Interpretación de lecturas actuales}

Las tarjetas principales muestran la última lectura disponible para la estación seleccionada.

\begin{table}[H]
\centering
\renewcommand{\arraystretch}{1.25}
\small
\begin{tabular}{|p{0.26\textwidth}|p{0.22\textwidth}|p{0.38\textwidth}|}
\hline
\rowcolor{gray!20}
\textbf{Variable} & \textbf{Unidad} & \textbf{Interpretación} \\
\hline
Humedad del suelo & \% & Variable principal para evaluar si el sustrato está seco, normal o húmedo. \\
\hline
Temperatura & $^\circ$C & Temperatura ambiente medida por la estación. \\
\hline
Humedad ambiental & \% & Humedad relativa del aire. \\
\hline
Luminosidad & lux & Intensidad luminosa recibida por el sensor BH1750. \\
\hline
Presión & hPa & Presión atmosférica medida por el sensor BME280. \\
\hline
\end{tabular}
\caption{Lecturas actuales mostradas por el \textit{dashboard}.}
\label{tab:lecturas_actuales_dashboard}
\end{table}

Si una variable aparece como ``---'', significa que no existe dato disponible para esa lectura, que el sensor no ha enviado valor o que todavía no se ha almacenado ninguna lectura válida.

\section{Estado de humedad del suelo}

El \textit{dashboard} clasifica la humedad del suelo en tres estados: seco, normal y húmedo. Esta clasificación se calcula utilizando los umbrales configurados en la API.

\begin{table}[H]
\centering
\renewcommand{\arraystretch}{1.25}
\small
\begin{tabular}{|p{0.24\textwidth}|p{0.28\textwidth}|p{0.34\textwidth}|}
\hline
\rowcolor{gray!20}
\textbf{Estado} & \textbf{Condición} & \textbf{Interpretación} \\
\hline
Seco & Humedad inferior al umbral bajo & El sustrato puede requerir riego según la lógica configurada en la estación. \\
\hline
Normal & Humedad entre umbral bajo y alto & El suelo se encuentra en una zona intermedia. \\
\hline
Húmedo & Humedad superior al umbral alto & El suelo dispone de humedad suficiente. \\
\hline
Sin dato & No hay lectura válida & No se debe interpretar el estado hídrico hasta recibir una lectura correcta. \\
\hline
\end{tabular}
\caption{Interpretación del estado de humedad del suelo.}
\label{tab:estado_humedad_suelo}
\end{table}

En la configuración por defecto, los umbrales utilizados por el \textit{dashboard} son 35\% para el límite bajo y 45\% para el límite alto. Estos valores pueden modificarse en el fichero de configuración \texttt{.env} mediante las variables \texttt{SOIL\_LOW\_THRESHOLD} y \texttt{SOIL\_HIGH\_THRESHOLD}.

\section{Estado de conexión y frescura de datos}

El \textit{dashboard} no solo muestra la última lectura, sino también su antigüedad. Esto permite saber si la estación está actualizando datos de forma reciente o si se ha perdido continuidad en la telemetría.

\begin{table}[H]
\centering
\renewcommand{\arraystretch}{1.25}
\small
\begin{tabular}{|p{0.24\textwidth}|p{0.30\textwidth}|p{0.32\textwidth}|}
\hline
\rowcolor{gray!20}
\textbf{Estado} & \textbf{Condición} & \textbf{Interpretación} \\
\hline
En vivo & Lectura reciente & La estación está enviando datos y la Raspberry Pi los está actualizando correctamente. \\
\hline
Datos recientes & Lectura algo antigua, pero dentro del margen permitido & El sistema tiene datos válidos, aunque conviene vigilar si continúa actualizando. \\
\hline
Sin actualización & Lectura antigua & Puede existir un problema de comunicación, publicación, ingesta o alimentación. \\
\hline
Sin datos & No existe ninguna lectura almacenada & La estación no ha publicado todavía o la Raspberry Pi no ha podido recuperar datos. \\
\hline
\end{tabular}
\caption{Estados de conexión y frescura de datos.}
\label{tab:estado_conexion_dashboard}
\end{table}

El tiempo utilizado para clasificar una lectura como reciente o antigua se controla mediante la variable \texttt{DATA\_STALE\_SECONDS}. En la configuración por defecto se utiliza un valor de 90 segundos.

\section{Estado del relé}

El \textit{dashboard} muestra el estado del relé recibido desde la estación. Este valor procede de la telemetría publicada en ThingSpeak y permite saber si el sistema de actuación está apagado o activo en el momento de la lectura.

\begin{table}[H]
\centering
\renewcommand{\arraystretch}{1.25}
\small
\begin{tabular}{|p{0.20\textwidth}|p{0.28\textwidth}|p{0.38\textwidth}|}
\hline
\rowcolor{gray!20}
\textbf{Valor} & \textbf{Estado mostrado} & \textbf{Interpretación} \\
\hline
0 & OFF & El relé está apagado y la bomba no debería estar activa. \\
\hline
1 & Activo & El relé está activo y el sistema de riego puede estar funcionando. \\
\hline
Sin dato & Sin dato & No se ha recibido información válida del relé. \\
\hline
\end{tabular}
\caption{Interpretación del estado del relé.}
\label{tab:estado_rele_dashboard}
\end{table}

El usuario debe tener en cuenta que el \textit{dashboard} muestra el estado recibido en la última lectura almacenada. Si existe retraso de comunicación, el estado físico real puede haber cambiado antes de que el \textit{dashboard} se actualice.

\section{Eventos de sistema}

El campo de evento permite interpretar el estado interno de la estación. La Tabla \ref{tab:eventos_sistema_usuario} recoge los códigos utilizados por el \textit{dashboard}.

\begin{table}[H]
\centering
\renewcommand{\arraystretch}{1.25}
\small
\begin{tabular}{|c|p{0.28\textwidth}|p{0.42\textwidth}|}
\hline
\rowcolor{gray!20}
\textbf{Código} & \textbf{Evento} & \textbf{Interpretación} \\
\hline
100 & \texttt{idle} & Estación en reposo, sin riego activo. \\
\hline
206 & \texttt{safe mode} & La estación trabaja en modo seguro. \\
\hline
210 & \texttt{riego ON} & La estación ha activado el riego. \\
\hline
211 & \texttt{riego OFF} & La estación ha finalizado el riego. \\
\hline
400 & \texttt{error} & Existe un fallo relevante o una lectura no válida. \\
\hline
429 & \texttt{cooldown} & El riego está bloqueado temporalmente por periodo de seguridad. \\
\hline
\end{tabular}
\caption{Eventos de sistema mostrados en el \textit{dashboard}.}
\label{tab:eventos_sistema_usuario}
\end{table}

Estos códigos ayudan a diferenciar una situación normal de reposo de un bloqueo por seguridad o de un error real del sistema.

\section{Códigos de diagnóstico}

El campo de diagnóstico aporta información auxiliar sobre el estado del ciclo de funcionamiento de la estación.

\begin{table}[H]
\centering
\renewcommand{\arraystretch}{1.25}
\small
\begin{tabular}{|c|p{0.28\textwidth}|p{0.42\textwidth}|}
\hline
\rowcolor{gray!20}
\textbf{Código} & \textbf{Etiqueta} & \textbf{Interpretación} \\
\hline
200 & OK & Funcionamiento correcto. \\
\hline
204 & no riego & No se ha activado el riego en el ciclo actual. \\
\hline
206 & modo seguro & La estación se encuentra en modo seguro. \\
\hline
400 & error & Se ha detectado un error. \\
\hline
429 & \textit{cooldown} & La estación evita un nuevo riego por seguridad temporal. \\
\hline
\end{tabular}
\caption{Códigos de diagnóstico mostrados en el \textit{dashboard}.}
\label{tab:codigos_diagnostico_usuario}
\end{table}

\section{Gráficas}

El \textit{dashboard} incluye gráficas de evolución reciente para las variables almacenadas. Estas gráficas permiten detectar tendencias de forma rápida, por ejemplo:

\begin{itemize}
    \item descenso progresivo de humedad del suelo;
    \item aumento de temperatura durante el día;
    \item variación de luminosidad;
    \item cambios de humedad ambiental;
    \item activaciones puntuales del relé;
    \item repetición de eventos de error o \textit{cooldown}.
\end{itemize}

Las gráficas se construyen a partir de las últimas lecturas almacenadas en SQLite. Si no existen datos suficientes, la gráfica mostrará un mensaje indicando que todavía no hay información disponible.

\section{Resumen histórico}

El resumen histórico muestra estadísticas de las últimas 24 horas para la estación seleccionada. Incluye número de muestras, fecha inicial y final de la ventana temporal, valores medios, mínimos y máximos.

\begin{table}[H]
\centering
\renewcommand{\arraystretch}{1.25}
\small
\begin{tabular}{|p{0.30\textwidth}|p{0.55\textwidth}|}
\hline
\rowcolor{gray!20}
\textbf{Dato histórico} & \textbf{Utilidad} \\
\hline
Número de muestras & Permite comprobar si la estación está enviando datos con regularidad. \\
\hline
Valor medio & Resume el comportamiento de una variable durante la ventana analizada. \\
\hline
Valor mínimo & Ayuda a detectar extremos, por ejemplo humedad mínima del suelo. \\
\hline
Valor máximo & Ayuda a detectar picos de temperatura, luminosidad o presión. \\
\hline
Primer y último registro & Permiten comprobar el intervalo real cubierto por los datos. \\
\hline
\end{tabular}
\caption{Interpretación del resumen histórico.}
\label{tab:resumen_historico_usuario}
\end{table}

\section{Tabla de últimas lecturas}

La tabla de últimas lecturas muestra registros recientes almacenados en la base de datos local. Cada fila corresponde a una lectura recuperada desde ThingSpeak y guardada por la Raspberry Pi.

La tabla permite revisar:

\begin{itemize}
    \item fecha y hora de cada lectura;
    \item humedad del suelo;
    \item temperatura;
    \item humedad ambiental;
    \item luminosidad;
    \item presión atmosférica.
\end{itemize}

Esta vista es útil para verificar si los datos evolucionan de forma coherente o si aparecen valores repetidos, ausentes o fuera de rango.

\section{Procedimiento de uso diario}

Una vez instalado el sistema, el uso habitual puede resumirse en los siguientes pasos:

\begin{enumerate}
    \item encender la estación remota y comprobar que queda alimentada correctamente;
    \item arrancar el sistema en Raspberry Pi mediante \texttt{./tierraviva.sh start};
    \item abrir el \textit{dashboard} desde un navegador;
    \item seleccionar la estación que se desea consultar;
    \item comprobar el estado de conexión;
    \item revisar la humedad del suelo;
    \item interpretar el estado del relé y los códigos de evento;
    \item consultar las gráficas si se desea analizar evolución;
    \item revisar la tabla de últimas lecturas ante cualquier valor anómalo;
    \item detener el sistema con \texttt{Ctrl+C} o mediante el comando de parada correspondiente si se está ejecutando como proceso gestionado.
\end{enumerate}

\section{Situaciones habituales y actuación recomendada}

\begin{table}[H]
\centering
\renewcommand{\arraystretch}{1.25}
\scriptsize
\begin{tabular}{|p{0.28\textwidth}|p{0.32\textwidth}|p{0.30\textwidth}|}
\hline
\rowcolor{gray!20}
\textbf{Situación observada} & \textbf{Interpretación probable} & \textbf{Actuación recomendada} \\
\hline
El \textit{dashboard} muestra ``Sin datos'' & No hay lecturas almacenadas para la estación seleccionada. & Comprobar que la estación publica en ThingSpeak y que \texttt{main.py} está ejecutándose. \\
\hline
El estado aparece como ``Sin actualización'' & La última lectura es demasiado antigua. & Revisar alimentación de la estación, cobertura LTE, canal ThingSpeak y \textit{logs} de Raspberry Pi. \\
\hline
La humedad aparece como ``Seco'' & El valor está por debajo del umbral bajo. & Verificar si la estación ejecuta riego local o si se mantiene bloqueada por \textit{cooldown}. \\
\hline
Aparece código 429 & La estación está en periodo de seguridad entre riegos. & Esperar a que finalice el \textit{cooldown} y revisar si el valor de humedad sigue siendo bajo. \\
\hline
Aparece código 400 & Se ha detectado un error o lectura no válida. & Revisar sensores, cableado, alimentación y monitor serie de la estación. \\
\hline
El relé aparece activo & La última telemetría indica relé encendido. & Comprobar físicamente si la bomba está funcionando y verificar que se apaga al finalizar el tiempo programado. \\
\hline
Las gráficas no muestran datos & No hay histórico suficiente o la API no devuelve lecturas. & Esperar nuevas lecturas y comprobar \texttt{/api/readings}. \\
\hline
Los valores no cambian durante mucho tiempo & Puede no haber nuevas publicaciones o el ingestor está ignorando duplicados. & Revisar entry\_id en ThingSpeak, \textit{logs} de \texttt{main.py} y estado de la estación. \\
\hline
\end{tabular}
\caption{Situaciones habituales durante el uso de TierraViva.}
\label{tab:situaciones_usuario_tierraviva}
\end{table}

\section{Buenas prácticas de uso}

Para utilizar TierraViva de forma segura y fiable se recomienda:

\begin{itemize}
    \item verificar que la estación arranca con el relé apagado;
    \item no dejar la bomba conectada durante pruebas iniciales sin supervisión;
    \item revisar periódicamente el estado de conexión del \textit{dashboard};
    \item comprobar que la humedad del suelo cambia de forma coherente;
    \item no modificar umbrales sin registrar el cambio;
    \item no compartir capturas que contengan claves API;
    \item mantener la electrónica protegida frente a agua y salpicaduras;
    \item revisar \textit{logs} si aparecen estados de error o datos obsoletos;
    \item contrastar el \textit{dashboard} con una observación física durante pruebas de riego.
\end{itemize}

\section{Conclusión}

El \textit{dashboard} de TierraViva permite supervisar de forma sencilla el estado de una o varias estaciones de riego inteligente. La interfaz resume las variables principales, clasifica el estado de humedad, muestra la frescura de los datos, representa gráficas recientes y permite revisar histórico almacenado en la Raspberry Pi.

El usuario debe interpretar el sistema como una herramienta de monitorización y apoyo a la validación del prototipo. La actuación sobre el riego se realiza en la estación Arduino mediante lógica local, mientras que el \textit{dashboard} permite comprobar si las lecturas, eventos y estados publicados son coherentes con el funcionamiento esperado.